\documentclass[11pt,a4paper]{article}
\usepackage[margin=25mm]{geometry}
\usepackage{fontspec}
\setmainfont{TeX Gyre Pagella}
\usepackage{amsmath,amssymb,mathtools}
\usepackage{graphicx}
\usepackage{xcolor}
\usepackage[unicode,colorlinks=true,linkcolor=blue,urlcolor=blue]{hyperref}
\usepackage{bookmark}
\usepackage{enumitem}
\setlist{itemsep=3pt,topsep=5pt}
\setlength{\emergencystretch}{3em}
\newcommand{\tightlist}{\setlength{\itemsep}{0pt}\setlength{\parskip}{0pt}}
\newcommand{\passthrough}[1]{#1}
\newcommand{\pandocbounded}[1]{#1}
\makeatletter
\def\maxwidth{\ifdim\Gin@nat@width>\linewidth\linewidth\else\Gin@nat@width\fi}
\def\maxheight{\ifdim\Gin@nat@height>0.75\textheight0.75\textheight\else\Gin@nat@height\fi}
\makeatother
\setkeys{Gin}{width=\maxwidth,height=\maxheight,keepaspectratio}
\hypersetup{pdfauthor={OpenAI GPT-6.1 Sol, Codex, Ultra effort}}
\begin{document}
\section{The norm index bound and Hasse's norm
theorem}\label{the-norm-index-bound-and-hasses-norm-theorem}

\emph{Written by OpenAI GPT-6.1 Sol in Codex, Ultra effort, October
2026. Self-checked by the writing AI; no independent review is claimed.
Public domain (CC0).}

The preceding lesson proved that a cyclic extension of degree \(n\) has
idèle-class norm index at least \(n\). We now prove the reverse bound.
Kummer theory supplies a subgroup of precisely the required index;
restriction and norm reduce arbitrary Galois groups to prime cyclic
extensions. In function-field characteristic \(p\), an adelic argument
supplies the case that roots of unity cannot reach.

The equality of the two bounds gives Hasse's norm theorem and the
local-global principle for conics. After the ray-prime theorem is proved
in lesson 22, its section 6A develops the full Hasse--Minkowski theorem
from these conic results.

We use \href{../hilberts-theorem-90-and-kummer-theory.html}{Hilbert's
Theorem 90 and Kummer theory},
\href{../ideles-in-extensions-and-their-cohomology.html}{Idèles in
extensions and their cohomology}, and
\href{../the-herbrand-quotient-of-the-idele-class-group.html}{The
Herbrand quotient of the idèle class group}. Write \(J_F\) for the
idèles, \(C_F=J_F/F^\times\), \(J_F^S\) for unrestricted components in
\(S\) and unit components elsewhere, and \(E_S=F^\times\cap J_F^S\).

\textbf{Prerequisite proof availability.} The named results below
identify specific programme lessons. The
\href{https://kokunoyumeto.github.io/open-math-courses/courses/NT-CFT/proof-dependencies.html\#lesson-15}{prerequisite
record} distinguishes published proofs, supplied owner texts awaiting
publication, and missing full proofs. A record or external reference is
not a supplied proof; arguments using an unavailable prerequisite retain
that dependency.

\subsection{1. The bound and the class field
axiom}\label{the-bound-and-the-class-field-axiom}

\textbf{Theorem 15.1.} If \(L/K\) is a finite Galois extension of global
fields, then \[
[C_K:N_{L/K}C_L]\leq [L:K].
\tag{1}
\] If the extension is cyclic of degree \(n\), then \[
\bigl|\widehat H^0(\operatorname{Gal}(L/K),C_L)\bigr|=n,
\qquad
\widehat H^{-1}(\operatorname{Gal}(L/K),C_L)=0.
\tag{2}
\] Thus idèle classes satisfy the cyclic class field axiom.

We first note that the quotient in (1) is finite, without using global
reciprocity. The idèle norm image is open, by the local norm theorem and
surjectivity on units outside finitely many places. Its image in \(C_K\)
is open. Also \[
|N_{L/K}x|_K=|x|_L.
\tag{3}
\] For number fields the right side takes every positive real value. For
function fields it takes a nonzero subgroup of the discrete group of
possible contents, so has finite index there. The compactness of
\(C_K^1\), proved in the preceding lesson, now makes the remaining
discrete quotient finite. Moreover this quotient is killed by \([L:K]\),
because the norm of an embedded \(K\)-idèle is its \([L:K]\)th power.

After proving (1), (2) follows immediately from the preceding lesson: \[
[C_K:NC_L]
=n\,\bigl|\widehat H^{-1}(G,C_L)\bigr|.
\tag{4}
\] The argument below proves (1) algebraically whenever the prime cyclic
degrees differ from the characteristic. Sections 5--7 prove the
remaining function-field case.

\subsection{2. Local power indices}\label{local-power-indices}

Let \(n=\ell^m\), with \(\ell\ne\operatorname{char}K\), and suppose
\(\mu_n\subset K\). For every place \(v\), \[
|K_v^\times/K_v^{\times n}|=\frac{n^2}{|n|_v}.
\tag{5}
\] At a finite place where \(n\) is a unit, deep principal units are
uniquely \(n\)-divisible: solve \((1+z)^n=1+t\) by successive residue
corrections, since the linear coefficient is a unit. Reduction
identifies the unit quotient with the residue-field power quotient, of
order \(n\); the \(n\) distinct roots of unity remain distinct in the
residue field. Valuations supply the second factor \(n\).

At a finite place dividing \(n\), put \(s_v=v(n)\). For \(k>s_v\), the
map \[
U^k\xrightarrow{(\cdot)^n}U^{k+s_v}
\tag{6}
\] is bijective. Indeed, on \(z\in\pi^k\mathcal O_v\), the equation \[
z=\frac{t-\sum_{j\geq2}\binom nj z^j}{n}
\] is a contraction when \(t\in\pi^{k+s_v}\mathcal O_v\): the nonlinear
part has Lipschitz valuation at least \(k-s_v>0\). The same estimate
proves injectivity. No nontrivial \(n\)th root of unity lies in \(U^k\).
The map \(U/U^k\to U^n/U^{k+s_v}\) therefore has kernel \(\mu_n\), of
order \(n\). Counting its image gives \[
[U:U^n]=nq_v^{s_v}.
\] Together with valuations this proves (5). In \(\mathbf C\) the
quotient is trivial and \(|n|_v=n^2\). A real place permits \(n>1\) only
when \(n=2\); its quotient has order \(2=4/|2|_v\).

Choose a finite, nonempty \(S\) containing the archimedean places, the
places dividing \(n\), and any specified ramified places, enlarged so
that \(J_K=K^\times J_K^S\). Such an \(S\) exists by the preceding
lesson. Put \(s=|S|\). The product formula applied to \(n\) gives \[
\left|\prod_{v\in S}K_v^\times/K_v^{\times n}\right|=n^{2s}.
\tag{7}
\] The \(S\)-unit theorem gives \[
E_S/E_S^n\simeq(\mathbf Z/n)^s.
\tag{8}
\] Here the rank is \(s-1\), and the finite cyclic torsion contains
\(\mu_n\). Also \(E_S/E_S^n\to K^\times/K^{\times n}\) is injective: an
\(n\)th root of an \(S\)-unit has zero valuations outside \(S\).

\subsection{3. The Kummer subgroup}\label{the-kummer-subgroup}

\textbf{Proposition 15.2.} Suppose \(\mu_n\subset K\), \(n=\ell^m\) is
prime to the characteristic, and \[
\operatorname{Gal}(L/K)\simeq(\mathbf Z/n)^r.
\] With \(S\) as above, including every ramified place, there is a set
\(T\) of \(s-r\) further finite places, all splitting completely in
\(L\), such that \[
\Delta=\{u\in E_S:u^{1/n}\in L\}
=\ker\left(E_S\longrightarrow
\prod_{v\in T}K_v^\times/K_v^{\times n}\right).
\tag{9}
\] For \[
I(S,T)=
\prod_{v\in S}K_v^{\times n}
\times\prod_{v\in T}K_v^\times
\times\prod_{v\notin S\cup T}U_v,
\qquad C(S,T)=\operatorname{im}(I(S,T)\to C_K),
\tag{10}
\] one has \[
[C_K:C(S,T)]=n^r=[L:K],
\qquad C(S,T)\subseteq N_{L/K}C_L.
\tag{11}
\]

\textbf{Proof.} Every Kummer radical class for \(L/K\) has an \(S\)-unit
representative. If \(a^{1/n}\in L\), then at an unramified place outside
\(S\), valuation in the unramified local extension shows \(n\mid v(a)\).
Form an idèle having valuation \(v(a)/n\) outside \(S\), with components
\(1\) in \(S\). Write it as \(b j\), with \(b\in K^\times\) and
\(j\in J_K^S\). Then \(a/b^n\in E_S\).

Let \(N=K(E_S^{1/n})\). By (8) and the perfect Kummer pairing proved in
lesson 3, \[
\operatorname{Gal}(N/K)\simeq(\mathbf Z/n)^s,\qquad L\subseteq N.
\] The restriction onto \(\operatorname{Gal}(L/K)\), a free
\(\mathbf Z/n\)-module, splits by lifting a basis. Its kernel \(H\) is
free of rank \(s-r\). For completeness, over \(\mathbf Z/\ell^m\), a
surjective matrix onto a free module has a unit pivot; elementary row
and column operations remove one pivot at a time, leaving a free kernel.
Choose a basis \(\sigma_1,\ldots,\sigma_{s-r}\) of \(H\).

Set \(F_i=N^{\langle\sigma_i\rangle}\). The cyclic prime-power extension
\(N/F_i\) has infinitely many unramified inert primes, by the stronger
splitting consequence proved in lesson 14. Choose one above a place
\(v_i\) of \(K\), avoiding \(S\) and previously chosen places. Its
decomposition group in the abelian extension \(N/K\) is cyclic of order
at most \(n\), and contains \(\langle\sigma_i\rangle\), which already
has order \(n\). Hence it equals \(\langle\sigma_i\rangle\). In
particular \(v_i\) splits in \(L\).

For \(u\in E_S\), its \(n\)th root is in \(K_{v_i}\) precisely when its
Kummer character kills \(\sigma_i\). Thus \(T=\{v_i\}\) gives (9). Each
\(v_i\) is outside the divisors of \(n\), and \(u\) is a unit there. The
kernel index is \(n^{s-r}\), which equals the order of
\(\prod_{v\in T}U_v/U_v^n\). Therefore \[
E_S\longrightarrow\prod_{v\in T}U_v/U_v^n
\quad\text{is surjective}.
\tag{12}
\] The case \(s=r\) simply has \(T=\varnothing\).

The crucial principal intersection is \[
K^\times\cap I(S,T)=E_{S\cup T}^{\,n}.
\tag{13}
\] One inclusion is clear. For the other, take \(y\) in the left side
and put \(M=K(y^{1/n})\), a cyclic extension of degree dividing \(n\).
Given any idèle class, represent it by \(x\in J_K^S\). Its components in
\(T\) are units. By (12), choose \(a\in E_S\) so that \(x/a\) is an
\(n\)th power at every place in \(T\). At places in \(S\), the extension
\(M\) splits, because \(y\) is already an \(n\)th power. Outside
\(S\cup T\), \(y\) is a unit and \(n\) is invertible, so \(M\) is
unramified and the unit \(x/a\) is a local norm. At \(T\), its being an
\(n\)th power makes it a local norm as well: every local degree divides
\(n\). Choose unit norm preimages outside finitely many places to obtain
an idèle norm. Thus \(NC_M=C_K\). The cyclic lower bound from lesson 14
forces \([M:K]=1\). Consequently \(y=b^n\), and its valuations show
\(b\in E_{S\cup T}\), proving (13).

Now \(J_K=K^\times J_K^{S\cup T}\), and \[
[J_K^{S\cup T}:I(S,T)]=n^{2s},\qquad
[E_{S\cup T}:E_{S\cup T}^{\,n}]=n^{s+|T|}=n^{2s-r}.
\] Dividing these indices using (13) proves the first assertion of (11).

Finally, every component in (10) is a local norm from \(L\). At \(S\),
the local abelian Galois group has exponent dividing \(n\); local
reciprocity identifies it with the local norm quotient, so that quotient
kills \(n\)th powers. This uses the exponent, not the potentially larger
local degree. At \(T\), the extension splits. Elsewhere it is unramified
and unit norms are surjective. Unit preimages at almost all places give
an idèle norm, proving the second assertion. \(\square\)

\subsection{4. Descent and arbitrary Galois
groups}\label{descent-and-arbitrary-galois-groups}

First let \(L/K\) be cyclic of prime degree
\(\ell\ne\operatorname{char}K\). Put \(K'=K(\mu_\ell)\), \(L'=LK'\), and
\(d=[K':K]\). Then \(d\mid\ell-1\), \(L\cap K'=K\), and \(L'/K'\) is
cyclic of degree \(\ell\). Proposition 15.2 gives index at most \(\ell\)
over \(K'\).

Restriction of idèles and norm back induce maps \[
C_K/NC_L\longrightarrow C_{K'}/NC_{L'}
\longrightarrow C_K/NC_L
\tag{14}
\] whose composite is \(x\mapsto x^d\). These maps are well-defined by
base change for field norms and transitivity; the same identities hold
componentwise for idèles and principal elements. The first quotient is
killed by \(\ell\), so the first map is injective. Its order is
therefore at most \(\ell\).

In any tower \(K\subset M\subset L\), \[
[C_K:N_{L/K}C_L]\leq
[C_K:N_{M/K}C_M]\,[C_M:N_{L/M}C_L].
\tag{15}
\] Indeed the kernel of the map between the first two quotients is the
image, under \(N_{M/K}\), of \(C_M/N_{L/M}C_L\). Prime towers prove (1)
for cyclic extensions and for Galois groups of prime-power order,
provided the relevant prime cyclic case is known.

Here is the reduction for every finite Galois group \(G\); a composition
series with cyclic factors would not suffice. Write \(|G|=p^a m\),
\(p\nmid m\). There is a subgroup \(P\) of order \(p^a\). An elementary
proof is to let \(G\) act on its subsets of size \(p^a\). The number
\(\binom{p^a m}{p^a}\) is congruent to \(m\) modulo \(p\), by taking the
coefficient of \(X^{p^a}\) in \((1+X)^{p^a m}\). Some orbit has size
prime to \(p\). Its stabilizer has order divisible by \(p^a\), and acts
freely on that subset, so its order also divides \(p^a\). This is the
required \(P\).

A nontrivial \(p\)-group has nontrivial center by its conjugacy-class
equation. Taking a central element of order \(p\) and inducting on the
quotient supplies a tower of prime cyclic extensions. Thus (15) bounds
the norm quotient for \(L/M\), where \(M=L^P\), by \(|P|\).

Restriction to \(M\) and norm back have composite
\(x\mapsto x^{[M:K]}\). Restriction is well-defined on the norm
quotients: on \(J_L\), group the product \(\prod_{g\in G}g(y)\) into the
cosets \(Pr\) to obtain a \(P\)-norm. Because \([M:K]\) is prime to
\(p\), restriction injects the \(p\)-primary part of \(C_K/NC_L\) into
\(C_M/N_{L/M}C_L\). Hence that primary part has order at most \(p^a\).
Multiplying over the primes dividing \(|G|\) proves (1).

This completes the algebraic argument for number fields and for prime
cyclic degrees different from the function-field characteristic. We next
supply that missing characteristic case without assuming global class
field theory.

\subsection{5. Additive duality for function
fields}\label{additive-duality-for-function-fields}

Let \(F\) have full constant field \(k=\mathbf F_q\) of characteristic
\(p\). The preceding lesson gave a separating parameter, a finite
separable extension \(F/E\) with \(E=k(t)\), and discreteness and
compact quotient for \(F\subset\mathbb A_F\).

On \(\mathbb A_E\), define \[
\Psi_E(x)=
\exp\left(\frac{2\pi i}{p}
\operatorname{Tr}_{k/\mathbf F_p}
\sum_v\operatorname{res}_v(x_v\,dt)\right).
\tag{16}
\] At a finite place the residue includes the residue-field trace to
\(k\). Only finitely many terms are nonzero. Partial fractions prove the
residue theorem for rational differentials on \(k(t)\): each finite
principal part contributes the negative of its contribution at infinity,
and a polynomial contributes zero. Thus \(\Psi_E\) is trivial on \(E\).

We need the stronger assertion \[
E^\perp=E
\quad\text{under }(x,y)\longmapsto\Psi_E(xy).
\tag{17}
\] Subtract from \(y\in E^\perp\) a rational function to put its
remainder in the compact representative set \[
\prod_{P\text{ finite}}\mathcal O_P\times t^{-1}k[[t^{-1}]]
\] constructed in lesson 14. Write its infinity component as
\(\sum_{j\geq1}c_jt^{-j}\). Testing against \(a t^{j-1}\), \(a\in k\),
gives \(\operatorname{Tr}_{k/\mathbf F_p}(a c_j)=0\), so all \(c_j\)
vanish. Now test against arbitrary proper fractions \(a(t)/P^r\). Other
places give zero, and these fractions represent every principal part at
\(P\). The local residue pairing between principal parts and
\(\mathcal O_P\) is nondegenerate: expand in the parameter \(P\), and
use nondegeneracy of the finite-field trace for the first nonzero
coefficient. The derivative \(P'(t)\) is a unit at \(P\), so replacing
\(dP\) by \(dt\) does not change this assertion. Each finite component
therefore vanishes. This proves (17).

The pairing also gives every continuous additive character of
\(\mathbb A_E\). Locally, a character is a linear map to the \(p\)th
roots of unity, killed by some parameter-power ideal. Coefficient
expansion and the finite-field trace identify it uniquely with a Laurent
series through the residue pairing. Globally, continuity kills all unit
components outside a finite set, so the representing Laurent series form
an adèle. Conversely an adèle gives a continuous character. This
establishes self-duality, not just nondegeneracy.

Put \(\Psi_F(x)=\Psi_E(\operatorname{Tr}_{F/E}x)\). Choose trace-dual
\(E\)-bases of \(F\). Under the topological isomorphism
\(\mathbb A_F\simeq\mathbb A_E^{[F:E]}\) proved in lesson 14, the
pairing \(\Psi_F(xy)\) is the coordinatewise pairing between those dual
bases. Thus \(\mathbb A_F\) is self-dual and \[
F^\perp=F.
\tag{18}
\] This argument avoids assuming a geometric residue or duality theorem
for \(F\).

For each \(w\), let \(N_w\) be the largest integer such that the local
character is trivial on \(\pi_w^{-N_w}\mathcal O_w\). It is nontrivial
at each place by the nondegenerate local trace. Outside finitely many
places the extension of \(E\) is unramified, its integral trace pairing
is perfect, and the base character has conductor \(N_w=0\). Hence \[
\kappa=\sum_w N_w w
\tag{19}
\] is a finite divisor. The local annihilator calculation from the
rank-one Fourier theory in lesson 12 gives, for every divisor
\(D=\sum d_w w\), \[
B(D)=\prod_w\pi_w^{-d_w}\mathcal O_w,\qquad
B(D)^\perp=B(\kappa-D).
\tag{20}
\]

\subsection{6. Riemann--Roch and finite-order Hecke
functions}\label{riemannroch-and-finite-order-hecke-functions}

Define \(H^0(D)=F\cap B(D)\), and \(l(D)=\dim_k H^0(D)\). It is finite:
a discrete subgroup meets a compact set in finitely many points. The
quotient \(\mathbb A_F/(F+B(D))\) is finite because \(\mathbb A_F/F\) is
compact and the image of \(B(D)\) is open. By (18)--(20), its characters
are exactly \(H^0(\kappa-D)\). A finite abelian group has as many
characters as elements, so that quotient has order \(q^{l(\kappa-D)}\).

Let \(c\) be the Haar volume of \(\mathbb A_F/F\). The image of \(B(D)\)
has volume \(\operatorname{vol}(B(D))/q^{l(D)}\). Counting its cosets
gives \[
\frac{\operatorname{vol}(B(D))}{c}
=q^{l(D)-l(\kappa-D)}.
\tag{21}
\] The ratio of the left sides for \(D\) and \(0\) is \(q^{\deg D}\).
Functions integral everywhere are the constants: they form a finite
\(k\)-algebra inside \(F\), hence a finite field contained in the full
constant field \(k\). Thus \(l(0)=1\). Put \(g=l(\kappa)\). Equation
(21) yields \[
l(D)-l(\kappa-D)=\deg D+1-g.
\tag{22}
\] Applying it to \(\kappa\) gives \(\deg\kappa=2g-2\). A
negative-degree divisor has no nonzero section, by the product formula.
Consequently \[
\deg D>2g-2\quad\Longrightarrow\quad l(D)=\deg D+1-g.
\tag{23}
\] This is the needed Riemann--Roch theorem, proved from the additive
quotient and its exact duality.

Now let \(\chi\) be a finite-order character of \(C_F\), with conductor
divisor \(\mathfrak f\). Let \(U(\mathfrak f)\) be deep principal units
at its support and full units elsewhere. The ray group \[
\Gamma_{\mathfrak f}=C_F/\operatorname{im}U(\mathfrak f)
\] is discrete, and its degree-zero subgroup is finite by compactness of
\(C_F^1\). Denote its order by \(h_{\mathfrak f}\), and write the degree
image as \(\delta\mathbf Z\). No assertion \(\delta=1\) is needed here.
Weak approximation makes any idèle principal times an idèle congruent to
\(1\) at \(\mathfrak f\); taking its other valuations shows that every
ray class is represented by a divisor prime to \(\mathfrak f\).

For \(\mathfrak f>0\), fix such a representative \(A\) of large degree
\(d\). The functions producing effective divisors in its ray class are
\[
a\in H^0(A),\qquad a\equiv1\pmod{\mathfrak f}.
\] Evaluation onto \(\mathcal O/\mathfrak f\) is surjective: its kernel
is \(H^0(A-\mathfrak f)\), and (23) makes the dimension difference
\(\deg\mathfrak f\), the dimension of the target. The indicated fiber
therefore has \[
q^{d-\deg\mathfrak f+1-g}
\tag{24}
\] elements. Different functions in this fiber give different divisors,
since the only constant congruent to \(1\) is \(1\). If
\(\mathfrak f=0\), divide the nonzero sections by \(k^\times\); the
count instead is \[
\frac{q^{d+1-g}-1}{q-1}.
\tag{25}
\] These counts depend only on the degree, not on the ray class.

The Euler product \[
L(\chi,T)=\prod_{w\notin\operatorname{supp}\mathfrak f}
(1-\chi(w)T^{\deg w})^{-1}
\tag{26}
\] is the generating function for these effective divisors with their
character weights. If \(\chi\) is nontrivial on the finite degree-zero
subgroup of \(\Gamma_{\mathfrak f}\), character summation on each
fixed-degree coset makes all sufficiently large coefficients zero. Hence
\(L(\chi,T)\) is a polynomial.

Otherwise \(\chi\) is a degree character, say
\(\chi(D)=\lambda^{\deg D}\), with \(|\lambda|=1\). It is unramified, so
its minimal conductor is zero. Formula (25) shows that \[
\zeta_F(T)
=\text{a polynomial}+
\sum_{\substack{d\geq d_0\\d\in\delta\mathbf Z}}
h_0\frac{q^{d+1-g}-1}{q-1}T^d.
\tag{27}
\] Thus its only possible denominator factors are \(1-T^\delta\) and
\(1-q^\delta T^\delta\). At \(T=q^{-1}\) it has a simple pole with
positive nonzero leading coefficient. A nontrivial degree character has
\(\lambda^\delta\ne1\), so \(L(\chi,T)=\zeta_F(\lambda T)\) is
holomorphic there. These divisor counts also justify absolute
convergence of the Euler products for \(|T|<q^{-1}\).

Replacing \(T\) by \(q^{-s}\), we have proved: \(\zeta_F(s)\) has a
simple pole at \(s=1\), and every nontrivial finite-order Hecke
\(L\)-function is holomorphic at \(1\). For number fields these exact
analytic inputs are the proved Theorems 10.1 and 10.2 and Corollary 10.4
in
\href{https://kokunoyumeto.github.io/open-math-courses/courses/NT-CFT/proof-dependencies.html\#NT-ADL-10}{Hecke
L-functions and the Dedekind zeta function}. The function-field proof
above uses neither reciprocity nor the norm index bound.

\subsection{7. Weber's argument}\label{webers-argument}

Let \(L/K\) be Galois of degree \(n\), and set \(A=C_K/NC_L\), a finite
group by section 1. Its characters are finite-order Hecke characters,
unramified outside a finite set \(S\). For \(v\notin S\), let
\(a_v\in A\) be the class of an idèle uniformizer. Character
orthogonality and the Euler logarithms give, for real \(s>1\), \[
\log\prod_{\chi\in A^\vee}L_S(s,\chi)
=|A|\sum_{\substack{v\notin S\\a_v=1}}(Nv)^{-s}+O(1).
\tag{28}
\] Indeed the terms of first degree have character sum \(|A|\) or \(0\);
all higher powers have bounded total near \(1\), since
\(\sum_v(Nv)^{-2s}\) converges there. The logarithm defined by the Euler
products is real, and their product is positive.

The principal factor has a simple pole. Each nonprincipal factor is
holomorphic at \(1\), by section 6 for function fields or the cited
internal analytic theorem for number fields. If its zero order there is
\(m_\chi\geq0\), taking real logarithms of absolute values in (28) gives
\[
|A|\sum_{a_v=1}(Nv)^{-s}
=\left(1-\sum_{\chi\ne1}m_\chi\right)
\log\frac1{s-1}+O(1).
\tag{29}
\] We have not assumed nonvanishing.

For the primes splitting completely in \(L\), the Euler logarithm of
\(\zeta_L\) gives \[
n\sum_{v\text{ split completely}}(Nv)^{-s}
=\log\frac1{s-1}+O(1).
\tag{30}
\] At a nonsplit unramified prime every residue degree is at least
\(2\), so its contribution is bounded by a constant times
\((Nv)^{-2s}\); ramified primes are finite. Every completely split prime
has \(a_v=1\), since its uniformizer is an idèle norm. Comparing (29)
and (30) proves \[
\frac{|A|}{n}\leq1-\sum_{\chi\ne1}m_\chi\leq1.
\] This proves (1) for all global fields, including characteristic-\(p\)
extensions of degree divisible by \(p\), and completes Theorem 15.1.

\subsection{8. Hasse's norm theorem and
conics}\label{hasses-norm-theorem-and-conics}

\textbf{Theorem 15.3 (Hasse's norm theorem).} For a finite cyclic
extension \(L/K\) of global fields, an element \(a\in K^\times\) is a
field norm if and only if it is a norm from \(L\otimes_K K_v\) at every
place \(v\).

\textbf{Proof.} Apply the cyclic Tate hexagon to \[
1\longrightarrow L^\times\longrightarrow J_L
\longrightarrow C_L\longrightarrow1.
\] Its relevant exact segment is \[
\widehat H^{-1}(G,C_L)\longrightarrow
K^\times/NL^\times\longrightarrow J_K/NJ_L.
\tag{31}
\] The first group is zero by (2). Thus the last map is injective. An
element in all local norm images gives an idèle norm: outside finitely
many ramified places it is a unit and has a unit norm preimage. Its
image in the last group is zero, so its field-norm class is zero. The
reverse implication follows by completion. \(\square\)

Cyclicity is essential; the theorem need not hold for general noncyclic
abelian extensions.

This cyclic result is also the input for the planned lesson \emph{The
principal genus theorem} in the Noether course, where the more general
cohomological comparison is developed.

\textbf{Corollary 15.4.} A smooth conic over a global field has a
rational point if and only if it has a point over every completion. In
characteristic different from \(2\), this applies to \[
aX^2+bY^2=Z^2,\qquad a,b\in K^\times,
\tag{32}
\] has a \(K\)-point if and only if it has a point over every \(K_v\).

\textbf{Proof.} If \(b\) is a square, \(X=0\) gives a point. Otherwise
put \(L=K(\sqrt b)\). A point in (32) must have \(X\ne0\), and the
equation becomes \[
a=N_{L/K}\left(Z/X+(Y/X)\sqrt b\right).
\] Over a completion where \(b\) becomes a square the split norm is
surjective; the conic has a point automatically. At other completions
the same equivalence holds. Theorem 15.3 proves the assertion.
\(\square\)

This treats every smooth conic in characteristic different from \(2\):
diagonalize its nondegenerate ternary form and scale the last
coefficient to obtain (32). Diagonal ternary equations in characteristic
\(2\) are not smooth conics. The corresponding smooth equation \[
X^2+XY+bY^2=aZ^2
\tag{33}
\] also has the local-global principle in characteristic \(2\): if
\(T^2+T+b\) splits there is an immediate point; otherwise the left side
is the norm of \(X+Y\theta\) from the cyclic Artin--Schreier quadratic
field. A point then has \(Z\ne0\), and Theorem 15.3 applies. Every
smooth characteristic-\(2\) conic either already has a rational point or
has this form. Its polar form has rank \(2\); choose a paired basis and
a radical vector to write its equation as \(A X^2+XY+B Y^2+C Z^2=0\).
Smoothness forces \(C\ne0\), since a zero of the quadratic form on its
polar radical would be a singular point. If \(A=0\), there is a rational
point. Otherwise divide by \(A\) and replace \(Y\) by \(A Y'\),
obtaining (33) with \(b=AB\) and \(a=C/A\). This completes the
characteristic-\(2\) case as well.

For example, in \(\mathbf Q(i)\) every odd prime \(p\equiv3\pmod4\) must
have even valuation in a norm. At \(2\), the norm subgroup is \[
2^{\mathbf Z}(1+4\mathbf Z_2),
\] as computed in the preceding idèle lesson. Thus \(3\) fails at both
\(3\) and \(2\), and \(7\) fails at both \(7\) and \(2\). The number
\(5=1^2+2^2\) is a norm globally.

\subsection{9. From conics to quadratic
forms}\label{from-conics-to-quadratic-forms}

The norm theorem proves the local-to-global principle for conics here.
The full Hasse--Minkowski theorem for nondegenerate quadratic forms over
number fields is proved in
\href{the-chebotarev-density-theorem.md\#6a-general-quadratic-forms}{The
Chebotarev density theorem, section 6A}, after reciprocity, existence
and the ray-prime theorem. That argument retains all dimensions and uses
the conic theorem above for its ternary base case.

\subsection{10. Exercises and complete
solutions}\label{exercises-and-complete-solutions}

\subsubsection{Exercise 1 --- Gaussian norms
(easy)}\label{exercise-1-gaussian-norms-easy}

Decide whether \(3,5,7\) are norms from \(\mathbf Q(i)\), identifying
local obstructions.

\textbf{Solution.} At an inert odd prime, norm valuations are even. The
valuations of \(3\) at \(3\), and \(7\) at \(7\), are odd, so both fail.
Their odd \(2\)-adic units are \(3\pmod4\), so they also fail at \(2\).
The element \(1+2i\) has norm \(5\); equivalently all its local tests
pass and Hasse's norm theorem applies. Positive sign at the real place
is satisfied by all three.

\subsubsection{Exercise 2 --- The cohomological injection
(medium)}\label{exercise-2-the-cohomological-injection-medium}

Derive Hasse's norm theorem directly from \(\widehat H^{-1}(G,C_L)=0\).

\textbf{Solution.} Invariants in \(L^\times\) and \(J_L\) are
\(K^\times\) and \(J_K\). The Tate hexagon for their exact sequence with
\(C_L\) gives (31). Vanishing makes the field-norm quotient inject into
the idèle-norm quotient. Local norm preimages can be chosen to be units
almost everywhere, so a locally norm element dies in the latter
quotient; injectivity kills its field-norm class. This explains both the
local-to-idèle step and the exact cohomological step.

\subsubsection{Exercise 3 --- The Kummer index
(medium)}\label{exercise-3-the-kummer-index-medium}

In Proposition 15.2, prove the exact index of \(C(S,T)\), including the
principal intersection.

\textbf{Solution.} The kernel of \(E_S\to\prod_T U_v/U_v^n\) is the
radical \(\Delta\), so its index is \(n^{s-r}\); this equals the target
order and proves surjectivity. For \(y\in K^\times\cap I(S,T)\), the
cyclic field \(K(y^{1/n})\) splits at \(S\), is unramified outside
\(S\cup T\), and has every idèle class as a norm after multiplying a
representative by the \(S\)-unit supplied by that surjectivity at \(T\).
The lower norm bound forces this field to be \(K\), and the root is an
\(S\cup T\)-unit. Hence the principal intersection is exactly
\(E_{S\cup T}^n\). The unrestricted \(S\)-local power quotients have
product order \(n^{2s}\) by (5) and the product formula; principal
\(S\cup T\)-units remove a subgroup of order \(n^{s+|T|}=n^{2s-r}\). The
resulting index is \(n^r\). Merely counting local factors without
proving the principal intersection would not establish it.

\subsubsection{Exercise 4 --- The analytic norm bound
(hard)}\label{exercise-4-the-analytic-norm-bound-hard}

Assume the simple pole of the global zeta functions at \(1\) and
holomorphy there of every nonprincipal finite-order Hecke function.
Prove (1) without assuming their nonvanishing.

\textbf{Solution.} The finite norm quotient \(A\) has \(|A|\)
characters. Orthogonality in its Euler products gives (28). If the
nonprincipal functions have total zero order \(m\), their product with
the principal one has logarithm \((1-m)\log(1/(s-1))+O(1)\). Thus the
primes with trivial uniformizer class have logarithmic density
\((1-m)/|A|\). Independently the Euler product of \(\zeta_L\) gives
density \(1/n\) for completely split primes: exactly those have \(n\)
terms of norm \(Nv\), and all other unramified terms have degree at
least \(2\). Complete splitting implies trivial norm-quotient class.
Therefore \[
\frac1n\leq\frac{1-m}{|A|}\leq\frac1{|A|},
\] proving \(|A|\leq n\). This also forces \(m=0\) for these characters.
Finite omitted Euler factors are nonzero at \(1\), so they change
neither zero orders nor the argument.

\subsection{Editable edition}\label{editable-edition}

The reading edition provides the complete LaTeX source of this lesson,
the cumulative course LaTeX and the editable source ZIP. The archive
contains all twenty-four Markdown lessons, complete LaTeX bodies,
original diagrams, metadata and reproduction instructions.

\subsection{References}\label{references}

The Kummer norm-index bound includes Sylow descent. The function-field
proof treats prime-to-characteristic radicals and characteristic-p
additive duality separately. The cyclic Hasse norm theorem and the
quadratic-form deduction retain their exact analytic and approximation
prerequisites.

\begin{itemize}
\tightlist
\item
  \href{https://www.jmilne.org/math/CourseNotes/CFT.pdf}{J. S. Milne,
  Class Field Theory, version 4.03}.
\item
  \href{https://kskedlaya.org/cft/sec_abstractcft1.html}{Kiran S.
  Kedlaya, Notes on class field theory, author-hosted HTML edition}.
\end{itemize}

The \href{../free-proofs.html}{proof guide} gives the lesson sequence
and the exact prerequisite record. External references accompany the
written arguments.

\end{document}
